\subsection{Function Signature Architecture}

A function signature is the call contract: how argument objects bind to local parameter names when the frame is created. Python supports positional and keyword parameters, defaults, positional-only (\texttt{/}), keyword-only (\texttt{*}), and annotations.

\subsubsection{Default Parameter Values and Definition-Time Evaluation}

Default values are evaluated when the \texttt{def} statement executes and stored on the function object (\texttt{\_\_defaults\_\_} as a tuple, \texttt{\_\_kwdefaults\_\_} as a dict for keyword-only defaults). Rebinding a name used in a default expression does not alter stored defaults.

%<<<PROGRAM:programs/fun02>>>


A function call in a default (\texttt{def f(x=make\_obj())}) runs once at definition time. Mutable defaults persist across calls---studied in the architectural pitfalls section.

\subsubsection{Positional-Only Parameters Using \texttt{/}}

Parameters before \texttt{/} accept positional arguments only. The marker is a binding-rule separator, not a parameter.

%<<<PROGRAM:programs/fun01>>>


\subsubsection{Keyword-Only Parameters Using \texttt{*}}

Parameters after a bare \texttt{*} must be supplied by keyword if provided at all. Required keyword-only parameters omit defaults.

\begin{verbatim}
def connect(host, *, port):
    return f"{host}:{port}"

print(connect("localhost", port=8080))
\end{verbatim}

Possible output:

\begin{verbatim}
localhost:8080
\end{verbatim}

\subsubsection{Function Annotations: Metadata for Tools, Not Automatic Runtime Type Enforcement}

Annotations attach metadata to parameters and return values in \texttt{\_\_annotations\_\_}. They do not enforce types at runtime; tools and linters may read them.

\begin{verbatim}
def build_line(name: str, score: int = 42) -> str:
    return f"{name}: {score}"

print(build_line("Homer"))
print(build_line("Homer", "not-an-int"))
print(build_line.__annotations__)
\end{verbatim}

Possible output:

\begin{verbatim}
Homer: Homer
Homer: not-an-int
{'name': <class 'str'>, 'score': <class 'int'>, 'return': <class 'str'>}
\end{verbatim}

A signature controls argument binding. Annotations describe intended meaning for readers and tools; they do not change the object-passing model.
